\chapter{04 - Ascend 310B4
掌纹识别工作台手动测试教程}\label{ascend-310b4-ux638cux7eb9ux8bc6ux522bux5de5ux4f5cux53f0ux624bux52a8ux6d4bux8bd5ux6559ux7a0b}

\section{案例目标}\label{ux6848ux4f8bux76eeux6807}

本案例在 Ascend 310B4 / 8T 开发板上运行 FastAPI + React
掌纹识别工作台。生产推理固定使用 NPU
\passthrough{\lstinline!mixed\_fp16!}；CPU、EDCC、origin
和离线候选比较不进入在线接口。

本版本是冻结的人工测试发布。CCNet
为默认/回滚模型；\passthrough{\lstinline!manual\_test!} profile
可逐一测试五个 CompNet。五个 CompNet 的状态为
\passthrough{\lstinline!manual\_test\_pending!}，不能把人工测试前的状态写成稳定验收通过。

\section{手动环境}\label{ux624bux52a8ux73afux5883}

在开发板同一个 shell 中执行：

\begin{lstlisting}[language=bash]
source /usr/local/miniconda3/etc/profile.d/conda.sh
conda activate base
source /usr/local/Ascend/ascend-toolkit/set_env.sh
cd <release-root>
export PALMPRINT_ROOT="$PWD"
export ACL_PYTHON_SITE=/usr/local/Ascend/ascend-toolkit/latest/python/site-packages
export PYTHONPATH="$PALMPRINT_ROOT:$ACL_PYTHON_SITE${PYTHONPATH:+:$PYTHONPATH}"
export PALMPRINT_PROFILE=manual_test
unset PALMPRINT_REQUIRE_TEMPLATE_ENCRYPTION
unset PALMPRINT_TEMPLATE_KEY_FILE
unset PALMPRINT_TEMPLATE_DIR
python -c "import sys, acl; print(sys.executable); print(acl.__file__)"
python -m palmprint_workbench.tools.verify_assets --strict
python -m palmprint_workbench.api --host 0.0.0.0 --port 7860
\end{lstlisting}

源码不会自动安装依赖、下载数据、构建 EDCC、转换模型或修改系统环境。OM
从固定来源取得后，必须按 \passthrough{\lstinline!om\_manifest.json!}
核对字节数和 SHA-256。

\section{API 和界面}\label{api-ux548cux754cux9762}

启动后检查
\passthrough{\lstinline!/api/health!}、\passthrough{\lstinline!/api/bootstrap!}、\passthrough{\lstinline!/api/candidates!}
和 \passthrough{\lstinline!/!}。manual\_test profile 的 bootstrap 应返回
CCNet 与五个 canonical CompNet，均为
\passthrough{\lstinline!npu/mixed\_fp16!}。

触摸屏逐页检查：系统状态、上传识别、模板注册、候选审计。测试
1280×720、1024×768 和 375×812；按钮不重叠，状态文本不溢出，API
不可用时显示空服务状态。

\section{六模型人工测试}\label{ux516dux6a21ux578bux4ebaux5de5ux6d4bux8bd5}

canonical ID 为：

\begin{itemize}
\tightlist
\item
  \passthrough{\lstinline!ccnet!}
\item
  \passthrough{\lstinline!compnet\_tongji\_600!}
\item
  \passthrough{\lstinline!compnet\_iitd\_460!}
\item
  \passthrough{\lstinline!compnet\_rest\_358!}
\item
  \passthrough{\lstinline!compnet\_xjtu\_flash\_200!}
\item
  \passthrough{\lstinline!compnet\_xjtu\_natural\_200!}
\end{itemize}

每个模型执行：上传合成 ROI
或匿名测试图、识别、注册三个模板、查询、命中识别、删除、重启后再查询。模板
namespace 固定为
\passthrough{\lstinline!<model\_id>\_\_npu\_\_mixed\_fp16!}。

每次记录模型 ID、OM SHA、请求时间、HTTP 状态、模型耗时、模板状态和服务
PID。当前内测版本会把用户主动触发的原图、ROI 和 JSON
索引自动保存到发行目录的
\passthrough{\lstinline!data/captures/!}，把明文
\passthrough{\lstinline!PWST1!} embedding 模板保存到
\passthrough{\lstinline!data/templates/!}；连续预览帧不保存。删除模板不会删除图片证据。真实掌纹图像、姓名、embedding
和完整报告不得进入源码或截图，也不得同步到 GitHub/HF。

查看和清理内测图片由操作者手动执行：

\begin{lstlisting}[language=bash]
find "$PALMPRINT_ROOT/data/captures" -type f
rm -rf "$PALMPRINT_ROOT/data/captures"
rm -rf "$PALMPRINT_ROOT/data/templates"
\end{lstlisting}

\section{摄像头测试}\label{ux6444ux50cfux5934ux6d4bux8bd5}

先检查 \passthrough{\lstinline!/api/cameras!}，优先选择 1280×720
MJPG。连续预览 100
帧，记录首帧时间、稳定帧间隔和实际输出尺寸；拍照识别后关闭预览，确认设备释放，再重新打开一次。若画面滞后，先降低预览分辨率并确认没有其他进程占用设备。

\section{常见问题}\label{ux5e38ux89c1ux95eeux9898}

\begin{longtable}[]{@{}
  >{\raggedright\arraybackslash}p{(\linewidth - 2\tabcolsep) * \real{0.5000}}
  >{\raggedright\arraybackslash}p{(\linewidth - 2\tabcolsep) * \real{0.5000}}@{}}
\toprule\noalign{}
\begin{minipage}[b]{\linewidth}\raggedright
现象
\end{minipage} & \begin{minipage}[b]{\linewidth}\raggedright
处理
\end{minipage} \\
\midrule\noalign{}
\endhead
\bottomrule\noalign{}
\endlastfoot
\passthrough{\lstinline!import acl!} 失败 & 重新加载同一 shell 的
conda/CANN，确认 Python 和 \passthrough{\lstinline!acl.\_\_file\_\_!}
来自同一环境 \\
模型不在列表 & 确认
\passthrough{\lstinline!PALMPRINT\_PROFILE=manual\_test!}、registry、OM
SHA 和 bootstrap \\
模板不存在 & 确认注册和识别使用同一模型、精度和当前发行目录的
\passthrough{\lstinline!data/templates!} \\
模板文件损坏 & 停止服务，备份
\passthrough{\lstinline!data/templates!}，再使用
\passthrough{\lstinline!.pstore.bak!}/\passthrough{\lstinline!.previous!}
恢复 \\
摄像头卡顿 & 检查 MJPG、实际分辨率和设备占用，先使用 960×540 或
1280×720 \\
服务健康但识别失败 & 执行真实上传或摄像头 smoke；health 不主动加载
NPU \\
ACL 异常退出、设备 reset 或进程意外退出 &
停止测试，保存时间/PID/模型/摘要，重启当前版本或回退
CCNet；不要升级版本 \\
\passthrough{\lstinline!Health: Alarm!} &
只记录设备状态，不单独判断模型或程序结果 \\
\end{longtable}

\section{版本冻结}\label{ux7248ux672cux51bbux7ed3}

人工测试期间不升级 Python、依赖、驱动、CANN
或应用版本。发现问题时停止当前模型、记录复现步骤、回退
CCNet，并在人工测试结束后单独规划下一版本。人工测试通过后，再更新准入证据、release
manifest、版本号和完整验收记录。
